\documentclass[11pt]{article}
\usepackage{preamble}
\setlength{\parskip}{4pt}

\begin{document}

\daybanner{AP Statistics \textperiodcentered\ Unit 3 \textperiodcentered\ Topic 3.7 \textperiodcentered\ B: 2026-12-07 \textperiodcentered\ E: 2027-01-13 \textperiodcentered\ TEACHER KEY}{Significant Does Not Mean Caused}

\sectionbanner{CED TETHER}

{\small
\begin{itemize}[leftmargin=1.5em,itemsep=2pt,topsep=2pt]
  \item \textbf{Skill 4.E} --- Justify a claim using a decision based on significance tests.
  \item \textbf{EU DAT-3} --- Significance testing allows us to make decisions about hypotheses within a particular context.
  \item \textbf{LO DAT-3.B} --- Justify a claim about the population based on the results of a significance test for a population proportion. [Skill 4.E]
  \item \textbf{EK DAT-3.B.1} --- The significance level, $\alpha$, is the predetermined probability of rejecting the null hypothesis given that it is true.
  \item \textbf{EK DAT-3.B.2} --- A formal decision explicitly compares the p-value to the significance level, $\alpha$. If the p-value $\leq$ $\alpha$, reject the null hypothesis. If the p-value > $\alpha$, fail to reject the null hypothesis.
  \item \textbf{EK DAT-3.B.3} --- Rejecting the null hypothesis means there is sufficient statistical evidence to support the alternative hypothesis. Failing to reject the null means there is insufficient statistical evidence to support the alternative hypothesis.
  \item \textbf{EK DAT-3.B.4} --- The conclusion about the alternative hypothesis must be stated in context.
  \item \textbf{EK DAT-3.B.5} --- A significance test can lead to rejecting or not rejecting the null hypothesis, but can never lead to concluding or proving that the null hypothesis is true. Lack of statistical evidence for the alternative hypothesis is not the same as evidence for the null hypothesis.
  \item \textbf{EK DAT-3.B.6} --- Small p-values indicate that the observed value of the test statistic would be unusual if the null hypothesis and probability model were true, and so provide evidence for the alternative. The lower the p-value, the more convincing the statistical evidence for the alternative hypothesis.
  \item \textbf{EK DAT-3.B.7} --- p-values that are not small indicate that the observed value of the test statistic would not be unusual if the null hypothesis and probability model were true, so do not provide convincing statistical evidence for the alternative hypothesis nor do they provide evidence that the null hypothesis is true.
  \item \textbf{EK DAT-3.B.8} --- A formal decision explicitly compares the p-value to the significance $\alpha$. If the p-value $\leq$ $\alpha$, then reject the null hypothesis, $H_0$: p = $p_0$. If the p-value > $\alpha$, then fail to reject the null hypothesis.
  \item \textbf{EK DAT-3.B.9} --- The results of a significance test for a population proportion can serve as the statistical reasoning to support the answer to a research question about the population that was sampled.
\end{itemize}
}
\textbf{Essential question:} How do the p-value, practical importance, and study design jointly limit a test conclusion?\par
\textbf{DOK-3 skill:} 4.E \textperiodcentered\ \textbf{FRQ pattern:} {\small\texttt{statistical-\allowbreak{}decision-\allowbreak{}vs-\allowbreak{}practical-\allowbreak{}size-\allowbreak{}and-\allowbreak{}scope}}\par
\textbf{DOK rationale:} Part (c) coordinates the formal test decision, conditional meaning of the p-value, practical importance, causal limits, and sampled-population scope to repair two partly defensible but overstated conclusions.\par

\frameworkphaseheader{Do Now (first take) $\rightarrow$ Explore (video + follow-along) $\rightarrow$ Exit (finish + turn in)}{1 $\rightarrow$ 3}{45}{%
\begin{itemize}[leftmargin=1.5em,itemsep=2pt,topsep=2pt]
  \item Board slide up at the bell. Record one trusted clause from each statement.
  \item Begin the video at minute 5; return at minute 33 to separate decision, effect size, and scope.
  \item Collect at the door. Score part (c) only, E/P/I.
\end{itemize}
}{%
\begin{itemize}[leftmargin=1.5em,itemsep=2pt,topsep=2pt]
  \item Mark one defensible and one questionable clause.
  \item Complete the follow-along, finish (a)--(c), revise, and turn in the sheet.
\end{itemize}
}{%
\begin{itemize}[leftmargin=1.5em,itemsep=2pt,topsep=2pt]
  \item What does 0.0354 assume about the September sorting rate?
  \item Which numbers govern the decision and describe the effect?
  \item What population and causal comparison does the design support?
\end{itemize}
}{Solo teacher. Circulate ~5 pairs. Require separate decision, importance, and scope statements.}

\begin{callout}{\IconWarn\ WATCH FOR}{calloutred}
\begin{itemize}[leftmargin=1.5em,itemsep=2pt,topsep=2pt]
  \item Reporting 0.0354 as the probability that the director's claim is true.
  \item Letting a small practical gap reverse the predetermined significance-test rule.
  \item Generalizing a September sample to every month or attributing the result to the labels.
\end{itemize}
\end{callout}

\sectionbanner{SCORING PART (c) --- E / P / I}

\textbf{Expected elements}
\begin{itemize}[leftmargin=1.5em,itemsep=2pt,topsep=2pt]
  \item \textbf{makes-formal-decision} (required) --- Uses $0.0354<0.05$ to reject and conclude $p<0.60$ for September
  \item \textbf{interprets-p-value} (required) --- Interprets 0.0354 under the null rather than as $P(H_0)$
  \item \textbf{separates-significance-and-importance} (required) --- Separates the formal decision from practical importance
  \item \textbf{bounds-cause-and-scope} (required) --- Limits scope to September and withholds label causation
  \item \textbf{states-reader-consequences} (required) --- States a consequence of each overreach
\end{itemize}

{\small\renewcommand{\arraystretch}{1.25}
\noindent\begin{tabular}{|p{0.5in}|p{5.9in}|}
\hline
\textbf{E} & Makes and interprets the decision, separates importance, bounds cause and scope, and repairs both overreaches. \\ \hline
\textbf{P} & Makes the decision but weakens one interpretation, scope, importance, or consequence. \\ \hline
\textbf{I} & Uses importance to reverse the test, treats 0.0354 as $P(H_0)$, accepts the null, or claims causation. \\ \hline
\end{tabular}}

\clearpage
\focusbanner{TODAY'S PROBLEM --- annotated}

A cafeteria served 5{,}000 September trays after new sorting-bin labels were installed. The director claims at least $60\%$ are sorted correctly. An SRS of 250 trays finds 136 correct; the club tests at $\alpha=0.05$.\par\smallskip \textbf{Imani:} ``The $p$-value is 0.0354, so reject $H_0$. The labels caused sorting below 60 percent for every month.''\par \textbf{Drew:} ``The 5.6-point gap may be too small to matter, so it is not significant. Accept the director's claim.''\par
\begin{center}
\textbf{September study}\par\smallskip
\begin{tabular}{|l|c|c|c|}
\hline
\textbf{Source} & \textbf{Trays} & \textbf{Correct} & \textbf{$\alpha$} \\ \hline
\textbf{September} & 5,000 & claim: $\geq 60\%$ & -- \\ \hline
\textbf{SRS} & 250 & 136 & 0.05 \\ \hline
\end{tabular}
\end{center}
\begin{firsttakebox}{FIRST TAKE (5 min)}
Which parts of Imani's and Drew's statements do you trust? Mark one statistical claim and one scope or importance claim to revisit.\par
\answer{Not graded. Each statement mixes a defensible idea with an overreach.}
\end{firsttakebox}

\textit{Rules callout ``MAKE THE DECISION, THEN BOUND THE CLAIM (keep on the page)'' is printed on the student sheet.}\par\medskip

\dokbadge{1}~\textbf{(a)} Define the population proportion $p$, state $H_0$ and $H_a$, and calculate $\widehat p$ and its signed difference from 0.60.\par
\answer{Let $p$ be the proportion of all 5{,}000 lunch trays served in September that were sorted correctly. Test $H_0:p=0.60$ versus $H_a:p<0.60$. The sample proportion is $\widehat p=136/250=0.544$, which is $0.544-0.600=-0.056$, or 5.6 percentage points below the null value.}

\dokbadge{2}~\textbf{(b)} Verify the Random, 10 percent, and Large Counts conditions. Calculate the one-proportion $z$ statistic and lower-tail $p$-value.\par
\answer{Random is supported by the stated SRS. The 10 percent condition holds because $250<0.10(5{,}000)=500$. Under $H_0$, $np_0=250(0.60)=150$ and $n(1-p_0)=250(0.40)=100$, both at least 10. Thus $SE_0=\sqrt{0.60(0.40)/250}=0.0309839$ and $z=(0.544-0.60)/0.0309839=-1.8074$. The lower-tail normal probability is $P(Z\leq-1.8074)=0.0354$.}

\dokbadge{3}~\textbf{(c)} Adjudicate both statements and give the strongest conclusion at $\alpha=0.05$. Interpret 0.0354 conditionally, distinguish significance from practical importance, bound cause and scope, and give the reader consequence of each overreach.\par
\answer{Imani's formal decision is warranted: $0.0354<0.05$, so reject $H_0$; there is convincing evidence that fewer than 60 percent of the 5{,}000 September trays were sorted correctly. Under $p=0.60$, 0.0354 is the chance that an SRS of 250 gives $\widehat p\leq0.544$, not the chance $H_0$ is true. Drew is right that practical importance needs a separate criterion, but it does not reverse the test rule or prove the null. The SRS supports September population inference; without random assignment or a comparison, it supports neither label causation nor every-month scope. Imani's overreach suggests unsupported cause and scope; Drew's suggests importance determines significance and a null can be accepted as true.}

\begin{summaryexitbox}{EXIT}
One thing the video changed about my first take: \hrulefill
\end{summaryexitbox}

\end{document}
